% ===== PRÉAMBULE CANONIQUE QCM — PHYSIQUE =====
% Sert à compiler controle.tex ET à la revalidation automatique du JSON
% (valider_qcm.py). NE PAS MODIFIER : packages figés.
% La bibliothèque TikZ "babel" est indispensable (conflit babel-french/circuitikz).
\documentclass[11pt,a4paper]{article}
\usepackage[utf8]{inputenc}\usepackage[T1]{fontenc}\usepackage{lmodern}
\usepackage[french]{babel}
\usepackage{amsmath,amssymb,mathtools}\usepackage{icomma}
\usepackage{siunitx}\sisetup{locale=FR}
\usepackage{xcolor}\usepackage{colortbl}
\usepackage{tikz}\usepackage{pgfplots}\pgfplotsset{compat=1.18}
\usepackage[siunitx,europeanresistors]{circuitikz}
\usepackage{enumitem}\usepackage{array,booktabs}
\usepackage[a4paper,margin=2.2cm]{geometry}
\usetikzlibrary{arrows.meta,calc,angles,quotes,positioning,patterns,%
  backgrounds,shapes.geometric,decorations.pathmorphing,%
  decorations.markings,intersections,babel}
\newcommand{\vv}[1]{\vec{#1}}\newcommand{\ds}{\displaystyle}
\newcommand{\R}{\mathbb{R}}\newcommand{\N}{\mathbb{N}}\newcommand{\Z}{\mathbb{Z}}
% ==============================================
\begin{document}

% =====================================================================
% QCM CONCOURS — PHYSIQUE — FICHE 06 « ONDES »
% 10 questions niveau concours difficile, 4 propositions / 1 correcte.
% Ce fichier calque le rendu EXACT du validateur (valider_qcm.py) :
%   \textbf{id --- theme} + énoncé + enumerate[label=\Alph*)] +
%   « Bonne réponse » + explication + itemize « Pièges ».
% =====================================================================

% ------------------------------------------------------------------ Q01
\noindent\textbf{PHYS-F06-Q01 --- Grandeurs caractéristiques : longueur d'onde}\par
Une onde progressive se propage le long d'une corde élastique à la célérité $v = \qty{72}{\kilo\metre\per\hour}$. Le graphe ci-dessous donne l'élongation $y$ d'un point de la corde en fonction du temps. La longueur d'onde de cette onde vaut :
\begin{center}
\begin{tikzpicture}
\begin{axis}[width=10cm,height=4.5cm,axis lines=middle,
  xlabel={$t$~(s)},ylabel={$y$~(cm)},
  xlabel style={below right},ylabel style={above left},
  xmin=0,xmax=1.58,ymin=-3.8,ymax=3.8,
  xtick={0.25,0.5,0.75,1,1.25,1.5},ytick={-3,3},
  tick label style={font=\footnotesize},grid=both,grid style={black!12},samples=200]
\addplot[blue,line width=1.1pt,domain=0:1.5]{3*sin(720*x)};
\end{axis}
\end{tikzpicture}
\end{center}
\begin{enumerate}[label=\Alph*)]
\item \qty{36}{\metre}
\item \qty{5}{\metre}
\item \qty{10}{\metre}
\item \qty{20}{\metre}
\end{enumerate}
\textit{Bonne réponse : C}\par
On lit la période sur le graphe : $T = \qty{0.5}{\second}$. On convertit la célérité : $v = \qty{72}{\kilo\metre\per\hour} = \qty{20}{\metre\per\second}$. Puis $\lambda = v\,T = 20 \times 0{,}5 = \qty{10}{\metre}$.\par
\textit{Pièges :}
\begin{itemize}
\item A : garde $v = 72$ sans convertir les \si{\kilo\metre\per\hour} en \si{\metre\per\second} : $72\times0{,}5 = \qty{36}{\metre}$.
\item B : lit une demi-période ($T = \qty{0.25}{\second}$, de la crête au creux) : $20\times0{,}25 = \qty{5}{\metre}$.
\item D : lit deux périodes comme une seule ($T = \qty{1}{\second}$) : $20\times1 = \qty{20}{\metre}$.
\end{itemize}
\vspace{8pt}\hrule\vspace{8pt}

% ------------------------------------------------------------------ Q02
\noindent\textbf{PHYS-F06-Q02 --- Acoustique : écho et échographie}\par
Lors d'une échographie, une sonde émet une brève salve d'ultrasons (célérité dans les tissus $v = \qty{1500}{\metre\per\second}$) et détecte l'écho renvoyé par un organe $\qty{80}{\micro\second}$ plus tard. La profondeur de cet organe est :
\begin{enumerate}[label=\Alph*)]
\item \qty{6.0}{\centi\metre}
\item \qty{12}{\centi\metre}
\item \qty{3.0}{\centi\metre}
\item \qty{24}{\centi\metre}
\end{enumerate}
\textit{Bonne réponse : A}\par
L'onde fait l'aller-retour sonde\,$\to$\,organe\,$\to$\,sonde : elle parcourt $2d$ pendant $\Delta t$. Donc $d = \dfrac{v\,\Delta t}{2} = \dfrac{1500 \times 80\times10^{-6}}{2} = \qty{6.0}{\centi\metre}$.\par
\textit{Pièges :}
\begin{itemize}
\item B : oublie le facteur $\tfrac12$ (ne tient pas compte de l'aller-retour) : $d = v\,\Delta t = \qty{12}{\centi\metre}$.
\item C : applique le facteur $\tfrac12$ deux fois : $d = v\,\Delta t/4 = \qty{3.0}{\centi\metre}$.
\item D : multiplie par $2$ au lieu de diviser : $d = 2\,v\,\Delta t = \qty{24}{\centi\metre}$.
\end{itemize}
\vspace{8pt}\hrule\vspace{8pt}

% ------------------------------------------------------------------ Q03
\noindent\textbf{PHYS-F06-Q03 --- Changement de milieu}\par
Un son de fréquence $f = \qty{500}{\hertz}$ se propage dans l'air (célérité $\qty{340}{\metre\per\second}$) puis pénètre dans l'eau, où sa célérité devient $v' = \qty{1360}{\metre\per\second}$. Sa longueur d'onde dans l'eau vaut :
\begin{enumerate}[label=\Alph*)]
\item \qty{0.68}{\metre}
\item \qty{2.04}{\metre}
\item \qty{0.17}{\metre}
\item \qty{2.72}{\metre}
\end{enumerate}
\textit{Bonne réponse : D}\par
La fréquence est imposée par la source : elle ne change pas d'un milieu à l'autre. Seules $v$ et $\lambda$ dépendent du milieu. Dans l'eau : $\lambda' = \dfrac{v'}{f} = \dfrac{1360}{500} = \qty{2.72}{\metre}$.\par
\textit{Pièges :}
\begin{itemize}
\item A : croit que $\lambda$ ne change pas et garde la célérité de l'air : $\lambda = 340/500 = \qty{0.68}{\metre}$.
\item B : utilise la différence des célérités : $(v'-v)/f = 1020/500 = \qty{2.04}{\metre}$.
\item C : inverse le rapport des célérités : $\lambda_{\text{air}}\times v/v' = 0{,}68\times0{,}25 = \qty{0.17}{\metre}$.
\end{itemize}
\vspace{8pt}\hrule\vspace{8pt}

% ------------------------------------------------------------------ Q04
\noindent\textbf{PHYS-F06-Q04 --- Lecture des graphes $y(t)$ et $y(x)$}\par
Les deux graphes ci-dessous décrivent la même onde progressive : à gauche l'élongation d'un point en fonction du temps $y(t)$, à droite le profil de la corde en fonction de la position $y(x)$ à un instant donné. La célérité de l'onde vaut :
\begin{center}
\begin{minipage}{0.48\textwidth}\centering
\begin{tikzpicture}
\begin{axis}[width=6.4cm,height=3.9cm,axis lines=middle,
  xlabel={$t$~(s)},ylabel={$y$~(cm)},xlabel style={below right},ylabel style={above left},
  xmin=0,xmax=1.28,ymin=-2.6,ymax=2.6,xtick={0.2,0.4,0.6,0.8,1,1.2},ytick={-2,2},
  tick label style={font=\scriptsize},grid=both,grid style={black!12},samples=180]
\addplot[blue,line width=1pt,domain=0:1.2]{2*sin(900*x)};
\end{axis}
\end{tikzpicture}
\end{minipage}\hfill
\begin{minipage}{0.48\textwidth}\centering
\begin{tikzpicture}
\begin{axis}[width=6.4cm,height=3.9cm,axis lines=middle,
  xlabel={$x$~(m)},ylabel={$y$~(cm)},xlabel style={below right},ylabel style={above left},
  xmin=0,xmax=6.4,ymin=-2.6,ymax=2.6,xtick={1,2,3,4,5,6},ytick={-2,2},
  tick label style={font=\scriptsize},grid=both,grid style={black!12},samples=180]
\addplot[orange,line width=1pt,domain=0:6]{2*sin(180*x)};
\end{axis}
\end{tikzpicture}
\end{minipage}
\end{center}
\begin{enumerate}[label=\Alph*)]
\item \qty{0.8}{\metre\per\second}
\item \qty{10}{\metre\per\second}
\item \qty{5}{\metre\per\second}
\item \qty{0.2}{\metre\per\second}
\end{enumerate}
\textit{Bonne réponse : C}\par
Le graphe $y(t)$ donne la période $T = \qty{0.4}{\second}$ ; le graphe $y(x)$ donne la longueur d'onde $\lambda = \qty{2}{\metre}$. La célérité vaut $v = \dfrac{\lambda}{T} = \dfrac{2}{0{,}4} = \qty{5}{\metre\per\second}$.\par
\textit{Pièges :}
\begin{itemize}
\item A : calcule $v = \lambda\,T = 2\times0{,}4 = \qty{0.8}{\metre\per\second}$ (multiplie au lieu de diviser).
\item B : lit $T = \qty{0.2}{\second}$ (une demi-période) : $v = 2/0{,}2 = \qty{10}{\metre\per\second}$.
\item D : intervertit les deux graphes ($\lambda$ lu sur $y(t)$, $T$ sur $y(x)$) : $v = 0{,}4/2 = \qty{0.2}{\metre\per\second}$.
\end{itemize}
\vspace{8pt}\hrule\vspace{8pt}

% ------------------------------------------------------------------ Q05
\noindent\textbf{PHYS-F06-Q05 --- Onde sinusoïdale : vitesse d'un point}\par
Une onde sinusoïdale progressive a pour équation (unités SI) :
\[ y(x,t) = 0{,}05\,\sin\!\left(20\pi\,t - 4\pi\,x\right). \]
La vitesse maximale d'oscillation d'un point du milieu vaut :
\begin{enumerate}[label=\Alph*)]
\item \qty{3.14}{\metre\per\second}
\item \qty{5.00}{\metre\per\second}
\item \qty{0.50}{\metre\per\second}
\item \qty{0.63}{\metre\per\second}
\end{enumerate}
\textit{Bonne réponse : A}\par
On identifie $A = \qty{0.05}{\metre}$, $\omega = 20\pi~\si{\radian\per\second}$ et $k = 4\pi~\si{\radian\per\metre}$. La vitesse maximale d'un point est $v_{\max} = A\,\omega = 0{,}05\times20\pi \approx \qty{3.14}{\metre\per\second}$. À ne pas confondre avec la célérité de propagation $v = \omega/k = \qty{5}{\metre\per\second}$.\par
\textit{Pièges :}
\begin{itemize}
\item B : confond avec la célérité de propagation $v = \omega/k = 20\pi/4\pi = \qty{5.00}{\metre\per\second}$.
\item C : utilise $A\,f$ (confond $\omega$ et $f$, oublie le facteur $2\pi$) : $0{,}05\times10 = \qty{0.50}{\metre\per\second}$.
\item D : utilise $A\,k$ au lieu de $A\,\omega$ : $0{,}05\times4\pi \approx \qty{0.63}{\metre\per\second}$.
\end{itemize}
\vspace{8pt}\hrule\vspace{8pt}

% ------------------------------------------------------------------ Q06
\noindent\textbf{PHYS-F06-Q06 --- Déphasage entre deux points}\par
Une onde sinusoïdale de longueur d'onde $\lambda = \qty{0.80}{\metre}$ se propage le long d'une corde. Deux points M et P de la corde sont distants de $d = \qty{2.0}{\metre}$. Le déphasage $\Delta\varphi = \dfrac{2\pi d}{\lambda}$ entre ces deux points vaut :
\begin{enumerate}[label=\Alph*)]
\item $2{,}5\pi$~rad
\item $0{,}8\pi$~rad
\item $3{,}2\pi$~rad
\item $5\pi$~rad
\end{enumerate}
\textit{Bonne réponse : D}\par
$\Delta\varphi = \dfrac{2\pi d}{\lambda} = \dfrac{2\pi \times 2{,}0}{0{,}80} = 5\pi$~rad. Modulo $2\pi$, cela équivaut à $\pi$ : les deux points vibrent en opposition de phase.\par
\textit{Pièges :}
\begin{itemize}
\item A : utilise $\pi d/\lambda$ (oubli du facteur $2$) : $2{,}5\pi$~rad.
\item B : inverse le rapport, $2\pi\lambda/d = 2\pi\times0{,}4$ : $0{,}8\pi$~rad.
\item C : multiplie par $\lambda$ au lieu de diviser, $2\pi d\lambda$ : $3{,}2\pi$~rad.
\end{itemize}
\vspace{8pt}\hrule\vspace{8pt}

% ------------------------------------------------------------------ Q07
\noindent\textbf{PHYS-F06-Q07 --- Ondes stationnaires : harmoniques}\par
Une corde de longueur $L = \qty{0.50}{\metre}$, fixée à ses deux extrémités, est le siège d'ondes de célérité $c = \qty{300}{\metre\per\second}$. La fréquence de la 3\textsuperscript{e} harmonique vaut :
\begin{enumerate}[label=\Alph*)]
\item \qty{300}{\hertz}
\item \qty{900}{\hertz}
\item \qty{1800}{\hertz}
\item \qty{450}{\hertz}
\end{enumerate}
\textit{Bonne réponse : B}\par
Pour une corde fixée aux deux extrémités, $f_n = \dfrac{n\,c}{2L}$. Le fondamental vaut $f_1 = \dfrac{300}{2\times0{,}50} = \qty{300}{\hertz}$, donc la 3\textsuperscript{e} harmonique $f_3 = 3f_1 = \qty{900}{\hertz}$.\par
\textit{Pièges :}
\begin{itemize}
\item A : donne le fondamental $f_1 = \qty{300}{\hertz}$ (oublie le facteur $n = 3$).
\item C : utilise $3c/L$ (oublie le $2$ au dénominateur : $\lambda_n = L/n$ au lieu de $2L/n$) : \qty{1800}{\hertz}.
\item D : applique $3c/4L$, formule de la corde à extrémité \emph{libre} ($4L$) au lieu de $2L$ : \qty{450}{\hertz}.
\end{itemize}
\vspace{8pt}\hrule\vspace{8pt}

% ------------------------------------------------------------------ Q08
\noindent\textbf{PHYS-F06-Q08 --- Ondes stationnaires : nœuds et ventres}\par
Une corde de longueur $L = \qty{1.2}{\metre}$, fixée à ses deux extrémités, vibre en régime stationnaire selon le mode ci-dessous (3 ventres). La longueur d'onde de cette onde stationnaire vaut :
\begin{center}
\begin{tikzpicture}[scale=1]
\def\Lh{9}\def\Ah{0.7}
\draw[green!55!black,line width=3pt] (0,-0.95)--(0,0.95);
\draw[green!55!black,line width=3pt] (\Lh,-0.95)--(\Lh,0.95);
\draw[blue,line width=1pt,domain=0:9,samples=150] plot (\x,{\Ah*sin(60*\x)});
\draw[black,line width=1pt,domain=0:9,samples=150] plot (\x,{-\Ah*sin(60*\x)});
\foreach \k in {0,1,2,3}{\fill (\k*3,0) circle (2.4pt);}
\draw[<->,>=Stealth,dashed,black!70] (0,-1.35)--(\Lh,-1.35) node[midway,fill=white,inner sep=1pt]{$L$};
\end{tikzpicture}
\end{center}
\begin{enumerate}[label=\Alph*)]
\item \qty{0.40}{\metre}
\item \qty{0.60}{\metre}
\item \qty{0.80}{\metre}
\item \qty{1.60}{\metre}
\end{enumerate}
\textit{Bonne réponse : C}\par
Une corde fixée aux deux bouts contient un nombre entier de fuseaux : $L = n\,\dfrac{\lambda}{2}$, d'où $\lambda = \dfrac{2L}{n}$. Avec $n = 3$ ventres et $L = \qty{1.2}{\metre}$ : $\lambda = \dfrac{2\times1{,}2}{3} = \qty{0.80}{\metre}$.\par
\textit{Pièges :}
\begin{itemize}
\item A : utilise $L/n$ (oublie le facteur $2$ dans $\lambda_n = 2L/n$) : $1{,}2/3 = \qty{0.40}{\metre}$.
\item B : compte les nœuds ($4$) au lieu des ventres ($3$) : $2L/(n{+}1) = 2{,}4/4 = \qty{0.60}{\metre}$.
\item D : applique $\lambda = 4L/n$ (formule d'une extrémité libre) : $4{,}8/3 = \qty{1.60}{\metre}$.
\end{itemize}
\vspace{8pt}\hrule\vspace{8pt}

% ------------------------------------------------------------------ Q09
\noindent\textbf{PHYS-F06-Q09 --- Interférences de deux sources}\par
Deux haut-parleurs $S_1$ et $S_2$ identiques, émettant \emph{en phase} un son de fréquence $f = \qty{680}{\hertz}$ (célérité du son dans l'air $v = \qty{340}{\metre\per\second}$). Un microphone placé en M est tel que $S_1\text{M} = \qty{4.25}{\metre}$ et $S_2\text{M} = \qty{2.00}{\metre}$ (voir figure). Que capte le microphone en M ?
\begin{center}
\begin{tikzpicture}[scale=0.85]
\coordinate (S1) at (0,0);\coordinate (S2) at (0,2.3);\coordinate (M) at (6,3.3);
\fill (S1) circle (2.4pt);\fill (S2) circle (2.4pt);\fill (M) circle (2.4pt);
\node[left] at (S1) {$S_1$};\node[left] at (S2) {$S_2$};\node[right] at (M) {M (micro)};
\draw[black,line width=0.8pt] (S1)--(M) node[midway,below,sloped,font=\footnotesize]{$d_1$};
\draw[black,line width=0.8pt] (S2)--(M) node[midway,above,sloped,font=\footnotesize]{$d_2$};
\end{tikzpicture}
\end{center}
\begin{enumerate}[label=\Alph*)]
\item Un son d'intensité maximale (interférence constructive).
\item Un son d'intensité minimale (interférence destructive).
\item Un son d'intensité moyenne : deux haut-parleurs distincts ne peuvent pas interférer.
\item On ne peut pas conclure sans connaître la distance $S_1S_2$.
\end{enumerate}
\textit{Bonne réponse : B}\par
$\lambda = v/f = 340/680 = \qty{0.5}{\metre}$. La différence de marche vaut $\delta = |d_1 - d_2| = 4{,}25 - 2{,}00 = \qty{2.25}{\metre} = 4{,}5\,\lambda = \dfrac{9\lambda}{2}$ : c'est un multiple impair de $\lambda/2$, donc interférence destructive (son minimal).\par
\textit{Pièges :}
\begin{itemize}
\item A : arrondit $\delta/\lambda = 4{,}5$ à un entier (ou inverse les conditions) et conclut à tort à une interférence constructive.
\item C : méconnaît la cohérence : de même fréquence et en phase, les deux sources \emph{sont} cohérentes et interfèrent.
\item D : croit nécessaire la distance $S_1S_2$ alors que les deux longueurs de trajet $d_1$ et $d_2$ suffisent ($\delta = d_1 - d_2$).
\end{itemize}
\vspace{8pt}\hrule\vspace{8pt}

% ------------------------------------------------------------------ Q10
\noindent\textbf{PHYS-F06-Q10 --- Fentes de Young : interfrange}\par
Une lumière monochromatique éclaire deux fentes de Young distantes de $a = \qty{0.24}{\milli\metre}$. Sur un écran placé à $D = \qty{2.0}{\metre}$, la distance entre la 1\textsuperscript{re} et la 5\textsuperscript{e} frange brillante est de $\qty{20}{\milli\metre}$ (voir figure). La longueur d'onde de cette lumière vaut :
\begin{center}
\begin{tikzpicture}[scale=0.85]
\draw[black,line width=1pt] (0,-1.6)--(0,1.6);\node[left,font=\footnotesize] at (0,0){$S$};
\draw[black,line width=1.4pt] (2,-1.6)--(2,-0.5);
\draw[black,line width=1.4pt] (2,-0.3)--(2,0.3);
\draw[black,line width=1.4pt] (2,0.5)--(2,1.6);
\node[right,font=\footnotesize] at (2.05,0.4){$F_2$};\node[right,font=\footnotesize] at (2.05,-0.4){$F_1$};
\draw[<->,>=Stealth] (2.55,-0.4)--(2.55,0.4) node[midway,right,font=\footnotesize]{$a$};
\draw[black,line width=1.4pt] (6.4,-1.9)--(6.4,1.9);\node[right,font=\footnotesize] at (6.4,1.7){écran};
\draw[<->,>=Stealth] (2,-1.85)--(6.4,-1.85) node[midway,below,font=\footnotesize]{$D$};
\foreach \y in {-1,-0.5,0,0.5,1}{\fill[orange] (6.4,\y) circle (1.7pt);}
\end{tikzpicture}
\end{center}
\begin{enumerate}[label=\Alph*)]
\item \qty{600}{\nano\metre}
\item \qty{480}{\nano\metre}
\item \qty{1200}{\nano\metre}
\item \qty{2400}{\nano\metre}
\end{enumerate}
\textit{Bonne réponse : A}\par
Entre la 1\textsuperscript{re} et la 5\textsuperscript{e} frange brillante, il y a $5-1 = 4$ interfranges : $i = 20/4 = \qty{5}{\milli\metre}$. Alors $\lambda = \dfrac{i\,a}{D} = \dfrac{5\times10^{-3}\times0{,}24\times10^{-3}}{2{,}0} = \qty{600}{\nano\metre}$.\par
\textit{Pièges :}
\begin{itemize}
\item B : divise l'écart par le nombre de franges ($5$) au lieu du nombre d'intervalles ($4$) : $i = \qty{4}{\milli\metre}$, $\lambda = \qty{480}{\nano\metre}$.
\item C : ne compte que $2$ intervalles entre les franges extrêmes : $i = \qty{10}{\milli\metre}$, $\lambda = \qty{1200}{\nano\metre}$.
\item D : prend l'écart total ($\qty{20}{\milli\metre}$) comme interfrange, sans diviser : $\lambda = \qty{2400}{\nano\metre}$.
\end{itemize}
\vspace{8pt}\hrule\vspace{8pt}

\end{document}
